\chapter{Loss of Taste (Ageusia/Dysgeusia)}

\section*{Definition}
Ageusia is the complete loss of taste function, while dysgeusia is a distorted taste sensation. True taste loss is rare; most complaints of taste loss actually reflect olfactory dysfunction. Taste sensation involves detection of five basic qualities: sweet, sour, salty, bitter, and umami.

\section*{What Happens in Your Body}
Taste perception begins when chemical stimuli activate taste receptor cells within taste buds on the tongue, soft palate, pharynx, and epiglottis. Signals travel via cranial nerves VII (facial), IX (glossopharyngeal), and X (vagus) to the nucleus tractus solitarius and then to the thalamus and gustatory cortex. Damage at any level impairs taste. Olfactory input strongly influences flavor perception, which is why nasal congestion distorts taste.

\section*{Causes}
\begin{itemize}
  \item Upper respiratory tract infections and COVID-19 (viral-induced olfactory/taste dysfunction)
  \item Nasal congestion, sinusitis, allergic rhinitis
  \item Head trauma (olfactory nerve or tract damage)
  \item Aging (progressive loss of taste buds and olfactory neurons)
  \item Medications: ACE inhibitors (dysgeusia), metronidazole, carbimazole, chemotherapy agents
  \item Radiation therapy to head and neck
  \item Oral conditions: Dry mouth, oral thrush, lichen planus, glossitis
  \item Neurological disorders: Alzheimer disease, Parkinson disease, multiple sclerosis
  \item Nutritional deficiencies: Zinc deficiency, vitamin B12 deficiency
  \item Cigarette smoking
  \item Idiopathic (no identifiable cause)
\end{itemize}

\section*{When to See a Doctor}
See your doctor if:
\begin{itemize}
  \item Persistent taste loss affecting quality of life or nutritional intake
  \item Sudden onset of taste loss (especially with other neurological symptoms)
  \item Associated weight loss, poor appetite, or depression
  \item Recent head trauma or neurological symptoms
  \item Concurrent loss of smell (anosmia)
  \item Burning mouth or oral pain
\end{itemize}

\section*{Related Symptoms}
\begin{itemize}
  \item Anosmia (loss of smell)
  \item Nasal congestion
  \item Dry mouth
  \item Oral thrush
  \item Weight loss
  \item Decreased appetite
  \item Depression
  \item Burning mouth sensation
\end{itemize}

\section*{Self-Care / Home Management}
\begin{itemize}
  \item Maintain good oral hygiene including tongue brushing.
  \item Zinc lozenges or supplementation after medical evaluation.
  \item Use plastic utensils if metallic taste persists.
  \item Add strong seasonings (herbs, spices) to enhance flavor of foods.
  \item Avoid smoking and limit alcohol.
  \item Stay hydrated to maintain saliva production.
  \item Consult physician for medication review (may contribute).
\end{itemize}

\section*{How Your Doctor Will Evaluate You}
\begin{itemize}
  \item Detailed history: onset, duration, associated symptoms, medications, infections, head trauma.
  \item Physical examination: oral cavity, cranial nerve function, nasal patency.
  \item Psychophysical taste testing (taste strips or solutions).
  \item Olfactory testing (Sniffin Sticks, UPSIT).
  \item MRI brain if central cause suspected (tumor, stroke, multiple sclerosis).
  \item Nutritional labs: zinc, vitamin B12, ferritin levels.
\end{itemize}

\section*{Treatment Options}
\begin{itemize}
  \item Treat underlying cause: Antibiotics for sinusitis, antifungal for thrush, smoking cessation.
  \item Olfactory training for post-viral/idiopathic loss.
  \item Medication adjustment or substitution if drug-induced.
  \item Zinc supplementation if deficient.
  \item Speech and language therapy for coping strategies.
  \item Cognitive behavioral therapy for associated depression or anxiety.
\end{itemize}

\section*{References}

\begin{thebibliography}{99}
\bibitem{harrison-taste} Longo DL, et al. \emph{Harrison's Principles of Internal Medicine}. 21st ed. McGraw-Hill; 2022. Chapter 34: Taste Disorders.
\bibitem{uptodate-taste} Bartoshuk LM, et al. Taste and smell disorders. In: UpToDate, Post TW (Ed), UpToDate, Waltham, MA; 2025.
\bibitem{snyder} Snyder DJ, et al. Taste dysfunction: a clinical review. \emph{N Engl J Med}. 2023;388(9):822-834.
\bibitem{heiser} Heiser C, et al. Post-COVID-19 taste loss: pathophysiology and therapy. \emph{J Neurol}. 2024;271(2):678-690.
\bibitem{hummel-taste} Hummel T, et al. Clinical assessment of gustatory function. \emph{Chem Senses}. 2023;48:bjad005.
\bibitem{green} Green BG, et al. \emph{The Sense of Taste}. 3rd ed. CRC Press; 2024. Chapter 15: Clinical Taste Disorders.
\end{thebibliography}

\clearpage
